\chapter{The Context Engine}
\label{ch:context}

\section{The objective function}

Every agentic editor faces the same problem and almost none of them state it as an
optimisation. Given a token budget $B$, choose a set of context items $S$ that maximises the
probability the agent produces a correct action.

\[
\max_{S} \; P(\text{correct action} \mid S) \quad \text{subject to} \quad \sum_{i \in S} \text{tokens}(i) \le B
\]

The naive implementation ranks chunks by embedding similarity and fills the budget until it
overflows. This is wrong in three distinct ways. It ignores that some items are load bearing
regardless of similarity, such as the exact signature of the function being modified. It
ignores that many high similarity items are redundant with each other. And it pays full
token price for information that could be represented at a fraction of the cost without
losing decision value.

Xyra treats this as a compression problem with a value model. The engine is built from four
mechanisms: structural representation, tiered escalation, delta encoding, and a budget
solver. Around them sit accelerators that reduce latency and cost without changing the
selected set.

\section{Structural Context Protocol}

The Structural Context Protocol, abbreviated SCP, is the canonical representation of a source
file inside Xyra. It is not the file text. It is a projection of the file's parse tree into a
form that preserves everything an agent needs to reason about the file's interface and
dependencies, while discarding the bodies it does not need.

An SCP record for a file contains the module path and language, the ordered list of exported
and internal symbols with their full signatures, type declarations in full, documentation
comments attached to symbols, the import list resolved to concrete paths, the outbound call
edges by symbol, and a stable content hash per symbol body so a body can be requested later
by handle.

\begin{table}[h]
\centering
\small
\begin{tabularx}{\textwidth}{l r r X}
\toprule
\textbf{Representation} & \textbf{Tokens} & \textbf{Ratio} & \textbf{Preserves} \\
\midrule
Raw file text & 3200 & 1.00 & Everything \\
SCP level 2 & 980 & 0.31 & Signatures, types, docs, imports, call edges, body handles \\
SCP level 1 & 420 & 0.13 & Signatures, types, imports \\
SCP level 0 & 130 & 0.04 & Module path, exported symbol names, import list \\
\bottomrule
\end{tabularx}
\caption{Indicative compression profile for a typical eight hundred line source file. Ratios are
representative of the design target, and the harness in Chapter~\ref{ch:quality} measures the
real distribution per language.}
\end{table}

The critical property is that SCP is \req{lossless with respect to the interface}. An agent
reading SCP level 2 knows every symbol it can call, with exact parameter types and return
types. It does not know how those symbols are implemented, and in the majority of tasks it
does not need to. When it does, it asks.

\section{Tiered escalation}

Context is assembled as a ladder, and the agent climbs it only where it must. This inverts
the usual flow, in which the editor guesses what the agent needs before the agent has
reasoned about the task.

\begin{figure}[h]
\centering
\begin{tikzpicture}[
  tier/.style={draw=xyraline,fill=xyrapanel,rounded corners=2pt,minimum height=9mm,minimum width=112mm,font=\small},
  arr/.style={-{Stealth[length=2mm]},xyralime}
]
\node[tier] (t0) at (0,0) {\textbf{Tier 0} \; Repository skeleton: directory shape, module names, entry points};
\node[tier,below=4mm of t0] (t1) {\textbf{Tier 1} \; SCP level 1 for the impact closure of the task};
\node[tier,below=4mm of t1] (t2) {\textbf{Tier 2} \; SCP level 2 for the files being edited plus their direct callers};
\node[tier,below=4mm of t2] (t3) {\textbf{Tier 3} \; Full bodies, fetched by handle on explicit request};
\node[tier,below=4mm of t3] (t4) {\textbf{Tier 4} \; Runtime evidence: test output, traces, screenshots, logs};
\draw[arr] (t0.east) -- ++(6mm,0) |- (t1.east);
\draw[arr] (t1.east) -- ++(6mm,0) |- (t2.east);
\draw[arr] (t2.east) -- ++(6mm,0) |- (t3.east);
\draw[arr] (t3.east) -- ++(6mm,0) |- (t4.east);
\end{tikzpicture}
\caption{The context ladder. Escalation is driven by the agent's own tool calls, not by
the editor's guess.}
\end{figure}

Tier~0 and Tier~1 are pushed into the session automatically because they are cheap and
always relevant. Tier~2 is computed from the deterministic impact closure described in
Section~\ref{sec:structural}. Tier~3 and Tier~4 are pulled by the agent through tools, which
means the cost is paid only when the information is actually wanted.

\section{Structural retrieval}
\label{sec:structural}

Before any embedding is consulted, Xyra answers the question structurally.

\begin{description}[style=nextline,leftmargin=0pt]
  \item[Import graph.] Parsed from source, resolved to concrete files, stored as edges. Gives
  the exact set of files that would break if a module changes, by transitive distance.
  \item[Symbol table.] Every definition with its kind, location and visibility, so a symbol
  name resolves to a definition without similarity search.
  \item[Call edges.] Outbound calls per symbol, which produces a call closure for a change to
  a single function rather than a whole file.
  \item[Test mapping.] Which test files exercise which source files, derived from imports and
  from execution traces when available. This makes the verification loop targeted rather than
  running the whole suite each time.
  \item[Blame and history.] Which commits last touched the region, which gives an agent the
  reason a line exists before it deletes it.
\end{description}

Embedding search is consulted for a narrow class of questions: locating behaviour described
in prose when the developer does not know the symbol name. It is a fallback, and the
interface labels its results as such so the developer can see when the machine is guessing.

\section{Semantic delta encoding}

Within a mission, consecutive sessions overlap heavily. Re-sending the same skeleton on every
ticket wastes a large fraction of the budget. Xyra addresses this with a delta protocol.

The core maintains, per mission, a content addressed dictionary of context items already
transmitted. When assembling the next session, items already in the dictionary are referenced
by handle rather than repeated, and the session preamble carries only the delta: what changed
since the last transmission, plus the handles for everything unchanged.

\begin{table}[h]
\centering
\small
\begin{tabularx}{\textwidth}{l X r}
\toprule
\textbf{Element} & \textbf{Encoding} & \textbf{Cost} \\
\midrule
Unchanged file, previously sent & Handle reference with symbol list & Tens of tokens \\
Changed file & Delta against the previous SCP record & Proportional to the change \\
New file in scope & Full SCP record at the required tier & Full record cost \\
Removed from scope & Explicit eviction notice & Negligible \\
\bottomrule
\end{tabularx}
\caption{Delta encoding between sessions in a mission.}
\end{table}

Because each ticket runs in a fresh session, the dictionary is not a context window trick. It
is a compact resumption record that lets a new session reach the same understanding as its
predecessor for a fraction of the tokens.

\section{The budget solver}

Given candidate items with an estimated value and a measured token cost, selection is a
knapsack problem with dependencies and redundancy penalties.

\begin{itemize}
  \item \req{Value} combines structural proximity to the task target, recency of change,
  authorship of the failing test, and explicit developer pinning. Similarity contributes, but
  it never dominates a structural signal.
  \item \req{Cost} is measured with a real tokeniser for the target model, not estimated by
  character count. Estimating by characters is wrong by up to forty percent on code.
  \item \req{Redundancy} is penalised. Two chunks with near identical content contribute once.
  \item \req{Dependencies} are respected. Including a symbol implies including the types in
  its signature.
  \item \req{Floors} are enforced. The file being edited is always at Tier~2 regardless of the
  solver's opinion, because omitting it produces confident nonsense.
\end{itemize}

The solver runs greedily by value density with a repair pass, which is fast enough to run on
every assembly and close enough to optimal that the difference does not matter in practice.

\section{Accelerators}

Compression reduces cost. Accelerators reduce latency and repeated work. They do not change
which items are selected, which keeps them safe to enable by default.

\begin{longtable}{@{}p{0.28\textwidth} p{0.64\textwidth}@{}}
\toprule
\textbf{Accelerator} & \textbf{Mechanism} \\
\midrule
\endhead
Stable prefix ordering & Context is emitted in a canonical order with the invariant material
first, so provider side prefix caches hit across sessions instead of missing on reordering. \\
Speculative prefetch & While the agent is reasoning, the graph predicts the next files it
will request and warms their SCP records and bodies. \\
Incremental embedding & Content hashed per chunk, so only changed chunks are re-embedded.
A file rename costs nothing. \\
Quantised vector index & Embeddings stored quantised and memory mapped, which keeps a large
repository index resident without dominating memory. \\
Local reranking & A small cross encoder run locally reorders the candidate set before the
budget solver, which improves selection quality without a network call. \\
Parallel lens execution & Independent review lenses run concurrently rather than in
sequence, so an adversarial review costs one lens of latency rather than three. \\
Warm sandbox worktrees & A pool of prepared worktrees with dependencies already installed,
so verification does not pay the setup cost on every ticket. \\
Test selection & The test mapping restricts verification to the tests that can be affected,
with a full suite run at ticket boundaries and before completion. \\
Streaming to disk & Agent output is persisted as it arrives, so a crashed session is
resumable and a long transcript never lives only in memory. \\
Compression of artefacts & Transcripts, screenshots and traces are stored compressed and
content addressed, deduplicated across missions. \\
\bottomrule
\caption{Latency and cost accelerators.}
\end{longtable}

\section{Context expansion integrations}

Structural knowledge of the repository is necessary but not sufficient. A large class of
agent errors comes from information that simply is not in the code. These integrations bring
that information into the ladder as first class context, each behind an explicit capability
so a developer can see and control what leaves the machine.

\begin{longtable}{@{}p{0.26\textwidth} p{0.30\textwidth} p{0.36\textwidth}@{}}
\toprule
\textbf{Integration} & \textbf{Brings in} & \textbf{Answers the question} \\
\midrule
\endhead
Language server & Hover types, call hierarchy, references, definitions & What is the real type here, who calls this \\
Version history & Blame, commit messages, pull request bodies & Why does this line exist \\
Issue tracker & Linked issues, acceptance criteria & What was actually asked for \\
Documentation fetch & Current library documentation by version & What is the correct API today, not at training time \\
Runtime traces & Execution paths, values at failure & What actually happened, not what should have \\
Test results & Failure output, flake history & Which failures are real \\
Design assets & Exported mockups, design tokens & What is this supposed to look like \\
Schema and migrations & Database schema, pending migrations & What shape is the data \\
Observability & Error rates, slow queries, recent incidents & What is breaking in production \\
Dependency registry & Package sizes, versions, advisories & What does adding this cost \\
Fleet siblings & Contracts consumed by other repositories & Who else breaks if this changes \\
Terminal history & Recent commands and their output & What has the developer already tried \\
\bottomrule
\caption{Context expansion integrations and the question each one closes.}
\end{longtable}

\section{The assembly pipeline}

\begin{figure}[h]
\centering
\begin{tikzpicture}[
  st/.style={draw=xyraline,fill=xyrapanel,rounded corners=2pt,minimum height=8.5mm,minimum width=27mm,font=\scriptsize,align=center},
  arr/.style={-{Stealth[length=1.8mm]},xyragrey}
]
\node[st] (task) {Task and\\target};
\node[st,right=7mm of task] (struct) {Structural\\closure};
\node[st,right=7mm of struct] (cand) {Candidate\\set};
\node[st,right=7mm of cand] (rank) {Local\\rerank};
\node[st,below=8mm of rank] (solve) {Budget\\solver};
\node[st,left=7mm of solve] (delta) {Delta\\encode};
\node[st,left=7mm of delta] (emit) {Emit in\\stable order};
\node[st,left=7mm of emit] (sess) {Session\\preamble};

\draw[arr] (task) -- (struct);
\draw[arr] (struct) -- (cand);
\draw[arr] (cand) -- (rank);
\draw[arr] (rank) -- (solve);
\draw[arr] (solve) -- (delta);
\draw[arr] (delta) -- (emit);
\draw[arr] (emit) -- (sess);

\node[st,above=8mm of cand,draw=xyralime] (vec) {Vector\\fallback};
\draw[arr,dashed] (vec) -- (cand);
\end{tikzpicture}
\caption{Context assembly. Similarity search enters only as a labelled fallback.}
\end{figure}

\section{Measurement}

None of this is worth shipping unless it is measured. The context engine carries a permanent
evaluation harness with four metrics, reported per mission and visible in the Context
Inspector.

\begin{itemize}
  \item \req{Budget utilisation.} Tokens spent versus tokens available, and the fraction
  spent on each tier.
  \item \req{Compression ratio.} Tokens emitted versus tokens the naive representation would
  have required for the same item set.
  \item \req{Escalation rate.} How often the agent had to climb to Tier~3, which is the
  signal that the lower tiers are under specified.
  \item \req{First attempt success.} The fraction of tickets whose first session passes
  verification, which is the only end to end measure that matters.
\end{itemize}
